\section{Conclusion}
\label{sec:conclusion}

We presented TopoPRM, a framework for mathematical reasoning post-training that combines deterministic verifiable process rewards with reverse-KL reasoning distillation.
The process reward module extracts dependency DAGs from reasoning traces and computes topology and continuity rewards through deterministic programs, requiring no learned reward model or human annotation.
The distillation stage transfers graph-enriched teacher behaviour to compact chain-like student reasoning via process-aware filtering and a mode-seeking reverse-KL objective.

TopoPRM improves overall critique accuracy by $+13.4$ points over outcome-only GRPO while reducing mean response length by 11\%.
The distilled Qwen3-8B student achieves 82.5\% on GSM8K and 59.8\% on MATH-500, competitive with models of comparable scale.
The method remains compatible with standard GRPO pipelines and introduces only plugin-level changes.

We scope our claims explicitly: deterministic reward computation provides verifiable structural process supervision, but does not constitute complete semantic proof verification.
Extending the framework to domains beyond mathematics and replacing the rule-based DAG extractor with a lightweight learned parser remain promising directions for future work.
